pdfoutput=1
% Options for packages loaded elsewhere
\PassOptionsToPackage{unicode}{hyperref}
\PassOptionsToPackage{hyphens}{url}
\pdfoutput=1
\documentclass[
  12pt,
]{article}
\usepackage{xcolor}
\usepackage{amsmath,amssymb}
\setcounter{secnumdepth}{-\maxdimen} % remove section numbering
\usepackage{iftex}
\ifPDFTeX
  \usepackage[T1]{fontenc}
  \usepackage{textcomp} % provide euro and other symbols
\else % if luatex or xetex
  \usepackage{unicode-math} % this also loads fontspec
  \defaultfontfeatures{Scale=MatchLowercase}
  \defaultfontfeatures[\rmfamily]{Ligatures=TeX,Scale=1}
\fi
\usepackage{lmodern}
\ifPDFTeX\else
  % xetex/luatex font selection
\fi
% Use upquote if available, for straight quotes in verbatim environments
\IfFileExists{upquote.sty}{\usepackage{upquote}}{}
\IfFileExists{microtype.sty}{% use microtype if available
  \usepackage[]{microtype}
  \UseMicrotypeSet[protrusion]{basicmath} % disable protrusion for tt fonts
}{}
\makeatletter
\@ifundefined{KOMAClassName}{% if non-KOMA class
  \IfFileExists{parskip.sty}{%
    \usepackage{parskip}
  }{% else
    \setlength{\parindent}{0pt}
    \setlength{\parskip}{6pt plus 2pt minus 1pt}}
}{% if KOMA class
  \KOMAoptions{parskip=half}}
\makeatother
\usepackage{longtable,booktabs,array}
\usepackage{calc} % for calculating minipage widths
% Correct order of tables after \paragraph or \subparagraph
\usepackage{etoolbox}
\makeatletter
\patchcmd\longtable{\par}{\if@noskipsec\mbox{}\fi\par}{}{}
\makeatother
% Allow footnotes in longtable head/foot
\IfFileExists{footnotehyper.sty}{\usepackage{footnotehyper}}{\usepackage{footnote}}
\makesavenoteenv{longtable}
\setlength{\emergencystretch}{3em} % prevent overfull lines
\providecommand{\tightlist}{%
  \setlength{\itemsep}{0pt}\setlength{\parskip}{0pt}}
\usepackage{bookmark}
\IfFileExists{xurl.sty}{\usepackage{xurl}}{} % add URL line breaks if available
\urlstyle{same}
\hypersetup{
  hidelinks,
  pdfcreator={LaTeX via pandoc}}

\author{SCX}
\date{2026}

\begin{document}

\section{SCX Self-Evolution: Edge Cases and Failure
Modes}\label{scx-self-evolution-edge-cases-and-failure-modes}

\begin{quote}
\textbf{Version}: 2026-06-28 \textbar{} \textbf{Status}: Formal analysis
\textbar{} \textbf{Prerequisite}: Documents 01, 02, 06, 10, 11
\textbf{Purpose}: Characterize failure modes of the SCX self-evolution
system, deriving conditions under which convergence fails, degrades, or
produces pathological behavior. Four canonical failure modes are
analyzed in depth, with formal conditions, rates of degradation, and
mitigation strategies. \textbf{Notation note}: \(\mathcal{S}\) denotes
the state space (as in core Theorems 1--3). \(S_t\) denotes the
gatekeeper scoring function at time \(t\) (as in Spring self-evolution
documents). \(\Psi\) denotes the Lyapunov function (renamed from
\(\Phi\) to avoid collision with feature space \(\Phi\) in Theorem 2).
\(\beta_t\) denotes the gatekeeper update rate; \(\eta\) is reserved for
the global noise rate.
\end{quote}

\begin{center}\rule{0.5\linewidth}{0.5pt}\end{center}

\subsection{Table of Contents}\label{table-of-contents}

\begin{enumerate}
\def\labelenumi{\arabic{enumi}.}
\tightlist
\item
  \hyperref[1-failure-mode-1-premature-convergence-gatekeeper-freezing]{Failure
  Mode 1: Premature Convergence (Gatekeeper Freezing)}
\item
  \hyperref[2-failure-mode-2-backlog-problem-memory-overrun]{Failure
  Mode 2: Backlog Problem (Memory Overrun)}
\item
  \hyperref[3-failure-mode-3-calibration-breakdown-client-divergence]{Failure
  Mode 3: Calibration Breakdown (Client Divergence)}
\item
  \hyperref[4-failure-mode-4-adversarial-poisoning-gatekeeper-corruption]{Failure
  Mode 4: Adversarial Poisoning (Gatekeeper Corruption)}
\item
  \hyperref[5-interaction-effects-compound-failures]{Interaction
  Effects: Compound Failures}
\item
  \hyperref[6-diagnostic-tests-for-each-failure-mode]{Diagnostic Tests
  for Each Failure Mode}
\item
  \hyperref[7-summary-table]{Summary Table}
\end{enumerate}

\begin{center}\rule{0.5\linewidth}{0.5pt}\end{center}

\subsection{1. Failure Mode 1: Premature Convergence (Gatekeeper
Freezing)}\label{failure-mode-1-premature-convergence-gatekeeper-freezing}

\subsubsection{1.1 Phenomenon}\label{phenomenon}

The gatekeeper converges to a suboptimal fixed point
\(S^*_{\text{sub}}\) before the student has learned enough to provide
useful feedback. Once frozen, the system cannot escape because the
gatekeeper rejects new data that would contradict its (incorrect)
judgments.

\subsubsection{1.2 Formal Mechanism}\label{formal-mechanism}

\textbf{The learning rate decays too fast.} When \(\beta_t\) (the
gatekeeper update rate) decays faster than the student can learn,
gatekeeper plasticity is lost before the student provides corrective
signal.

\begin{quote}
\textbf{Notation note (2026-06-28, B4):} The symbol \(\eta\) is reserved
throughout the SCX theory for the global noise rate \(\eta = P(Z=1)\).
The gatekeeper update rate is denoted \(\beta_t\). This document
previously used \(\eta(t)\) locally for the gatekeeper update rate,
which has been corrected to \(\beta_t\) for consistency.
\end{quote}

\textbf{Definition 1.1 (Freezing Time).} The \textbf{freezing time}
\(t_{\text{freeze}}\) is the earliest time such that for all
\(t \geq t_{\text{freeze}}\):

\[\|\Delta S_t\|_{M_0} \leq \varepsilon_{\text{mach}},\]

i.e., subsequent gatekeeper updates are numerically indistinguishable
from zero.

\textbf{Theorem 12.1 (Freezing Condition --- PROVEN).} Under condition
C6' (two-timescale, \(\beta_t = o(\alpha_t)\)), if:

\[\beta_t \cdot B_S \leq \varepsilon_{\text{mach}} \quad \text{for } t \geq t_{\text{freeze}},\]

then the gatekeeper freezes. With \(\beta_t = \beta_0 t^{-b}\), this
occurs at:

\[\boxed{\;t_{\text{freeze}} \leq \left(\frac{\beta_0 B_S}{\varepsilon_{\text{mach}}}\right)^{1/b}\;}.\]

\emph{Proof.} By C7, \(\|\Delta S_t\|_\infty \leq \beta_t B_S\). When
this bound falls below machine precision, the update is numerically
zero. Solving \(\beta_0 t^{-b} B_S = \varepsilon_{\text{mach}}\) gives
the freezing time. \(\square\)

\textbf{Status: PROVEN.} This is a direct consequence of C7 and the
numerical precision constraint.

\subsubsection{1.3 When Freezing is
Premature}\label{when-freezing-is-premature}

Freezing is \textbf{premature} if
\(t_{\text{freeze}} < t_{\text{converge}}\), where
\(t_{\text{converge}}\) is the time needed for the student to reach its
asymptotic accuracy.

\textbf{Theorem 12.2 (Premature Freezing Condition --- PROVEN).} Under
strong convexity (SC), the student needs:

\[t_{\text{converge}} = \left(\frac{C_1}{\varepsilon_{\text{target}}}\right)^{1/a}\]

steps to achieve
\(\mathbb{E}[\|\theta_t - \theta^*\|^2] \leq \varepsilon_{\text{target}}\).
Premature freezing occurs when:

\[\boxed{\;\left(\frac{\beta_0 B_S}{\varepsilon_{\text{mach}}}\right)^{1/b} < \left(\frac{C_1}{\varepsilon_{\text{target}}}\right)^{1/a}\;}.\]

\emph{Numerical example.} With \(\beta_0 = 0.1\), \(B_S = 1\),
\(\varepsilon_{\text{mach}} = 10^{-16}\), \(b = 0.8\): -
\(t_{\text{freeze}} \approx (0.1 / 10^{-16})^{1.25} \approx 10^{18.75}\),
astronomically large. Freezing is \textbf{not} a practical concern with
double precision.

With \(\beta_0 = 0.01\), \(\varepsilon_{\text{mach}} = 10^{-4}\)
(aggressive early stopping threshold), \(b = 0.9\): -
\(t_{\text{freeze}} \approx (0.01 / 10^{-4})^{1.11} \approx 10^{2.22} \approx 166\),
which is small enough to be concerning.

\textbf{Status: PROVEN.} The existence of a freezing time is guaranteed;
whether it is premature depends on parameter choices.

\subsubsection{1.4 Effect of Premature Freezing on Solution
Quality}\label{effect-of-premature-freezing-on-solution-quality}

\textbf{Theorem 12.3 (Suboptimality Gap Under Premature Freezing ---
PROVEN).} If the gatekeeper freezes at \(S_{t_{\text{freeze}}}\) and the
student continues to converge toward
\(\theta^*(S_{t_{\text{freeze}}})\), the asymptotic Lyapunov gap is:

\[\Psi(S_{t_{\text{freeze}}}, \theta^*(S_{t_{\text{freeze}}})) - \Psi(S^*, \theta^*) \geq \lambda \cdot \bigl(\mathbb{E}_{V_0}[\ell(f_{\theta^*(S_{t_{\text{freeze}}})}, y)] - \mathbb{E}_{V_0}[\ell(f_{\theta^*}, y)]\bigr) \geq 0.\]

The gap is non-negative and equals zero only if
\(S_{t_{\text{freeze}}} = S^*\).

\emph{Proof.} At the frozen gatekeeper \(S_{t_{\text{freeze}}}\), the
student converges to \(\theta^*(S_{t_{\text{freeze}}})\), which
minimizes the loss under the frozen distribution
\(P_{S_{t_{\text{freeze}}}}\). The gap on the reference distribution
\(V_0\) is generally positive because the frozen gatekeeper may exclude
regions where the true fixed-point student \(\theta^*\) performs
differently. \(\square\)

\subsubsection{1.5 Rate of Degradation}\label{rate-of-degradation}

\textbf{Corollary 12.1 (Degradation Rate Under Freezing).} At time
\(t \geq t_{\text{freeze}}\), the convergence of \(\Psi_t\) stalls:

\[\Psi_t - \Psi^* = \Theta(1) \quad \text{(does not decay further)}.\]

The improvement after freezing is:

\[\Delta\Psi_{\text{post-freeze}} = \Psi_{t_{\text{freeze}}} - \lim_{t \to \infty} \Psi_t = O(\alpha_{t_{\text{freeze}}}),\]

which comes only from the student continuing to train on the frozen
distribution.

\subsubsection{1.6 Mitigation Strategies}\label{mitigation-strategies}

\begin{longtable}[]{@{}
  >{\raggedright\arraybackslash}p{(\linewidth - 4\tabcolsep) * \real{0.3704}}
  >{\raggedright\arraybackslash}p{(\linewidth - 4\tabcolsep) * \real{0.4074}}
  >{\raggedright\arraybackslash}p{(\linewidth - 4\tabcolsep) * \real{0.2222}}@{}}
\toprule\noalign{}
\begin{minipage}[b]{\linewidth}\raggedright
Strategy
\end{minipage} & \begin{minipage}[b]{\linewidth}\raggedright
Mechanism
\end{minipage} & \begin{minipage}[b]{\linewidth}\raggedright
Cost
\end{minipage} \\
\midrule\noalign{}
\endhead
\bottomrule\noalign{}
\endlastfoot
\textbf{Minimum learning rate} & Clamp \(\beta_t \geq \beta_{\min} > 0\)
& Prevents asymptotic convergence, introduces steady-state error
\(O(\beta_{\min})\) \\
\textbf{Cyclical learning rate} & Periodically reset \(\beta_t\) to
\(\beta_0\) & Requires tuning cycle period; may cause transient
instability \\
\textbf{Plateau detection + boost} & When \(\|\Delta S_t\| < \delta\)
for \(k\) consecutive steps, multiply \(\beta_t\) by \(\gamma > 1\) &
Ad-hoc; may overshoot \\
\textbf{Student-driven unfreezing} & When student loss on \(V_0\) stops
improving, increase \(\beta_t\) temporarily & Couples unfreezing to
observable metric \\
\textbf{Polyak averaging on gatekeeper} & Maintain
\(\bar{S}_t = \frac{1}{t}\sum_{i=1}^t S_i\) & Reduces variance; may
smooth past frozen states \\
\end{longtable}

\begin{center}\rule{0.5\linewidth}{0.5pt}\end{center}

\subsection{2. Failure Mode 2: Backlog Problem (Memory
Overrun)}\label{failure-mode-2-backlog-problem-memory-overrun}

\subsubsection{2.1 Phenomenon}\label{phenomenon-1}

The memory bank \(M_t\) grows faster than the gatekeeper can score
incoming samples. Unscored samples accumulate in a backlog, delaying the
gatekeeper's feedback and eventually stalling the self-evolution loop.

\subsubsection{2.2 Formal Mechanism}\label{formal-mechanism-1}

\textbf{Assumptions:} - Samples arrive at rate \(r_{\text{in}}\)
(samples per unit time). - The gatekeeper can score at rate
\(r_{\text{score}}\) (samples per unit time). - The NEP student retrains
every \(K\) new samples.

\textbf{Definition 2.1 (Backlog).} The backlog at time \(t\) is:

\[B(t) = r_{\text{in}} \cdot t - r_{\text{score}} \cdot t = (r_{\text{in}} - r_{\text{score}}) \cdot t.\]

A backlog develops when \(r_{\text{in}} > r_{\text{score}}\).

\subsubsection{2.3 Gatekeeper Scoring
Complexity}\label{gatekeeper-scoring-complexity}

The gatekeeper's per-sample scoring cost is:

\[T_{\text{score}}(x) = O(M) \quad \text{(evaluate all } M \text{ experts)} + O(K_S) \quad \text{(state lookup)} + O(d_\phi) \quad \text{(feature computation)}.\]

Total scoring throughput:

\[r_{\text{score}} = \frac{1}{T_{\text{score}}} = \Omega\!\left(\frac{1}{M + K_S + d_\phi}\right).\]

\subsubsection{2.4 Critical Condition for
Backlog}\label{critical-condition-for-backlog}

\textbf{Theorem 12.4 (Backlog Instability --- PROVEN).} If:

\[\boxed{\;r_{\text{in}} > \frac{C_{\text{compute}}}{M + K_S + d_\phi}\;},\]

where \(C_{\text{compute}}\) is the available compute (FLOPs/second),
the backlog grows without bound and the system becomes
\textbf{unstable}: the memory bank contains samples whose scores are
based on a stale gatekeeper \(S_{t - \tau}\) with delay
\(\tau = B(t) / r_{\text{score}} \to \infty\).

\emph{Proof.} The scoring delay per sample is
\(\tau = B(t) / r_{\text{score}}\). When
\(r_{\text{in}} > r_{\text{score}}\), \(B(t) \to \infty\) and
\(\tau \to \infty\). The gatekeeper used for sample \(x_i\) is
\(S_{t_i}\) where \(t_i = i - \tau_i\). As \(\tau \to \infty\), the
gatekeeper version used for scoring becomes arbitrarily old, violating
the implicit assumption that \(S_t\) reflects current knowledge.
\(\square\)

\textbf{Status: PROVEN.} This is a throughput calculation.

\subsubsection{2.5 Effect on Self-Evolution
Dynamics}\label{effect-on-self-evolution-dynamics}

\textbf{Theorem 12.5 (Stale Gatekeeper Error --- CONJECTURED).} If
sample \(x\) is scored by gatekeeper \(S_{t-\tau}\) (stale by \(\tau\)
steps), the probability of misclassification is:

\[\mathbb{P}(\text{sign}(S_{t-\tau}(x) - \gamma) \neq \text{sign}(S_t(x) - \gamma)) \leq \frac{L_S \cdot \beta_{\max} \cdot \tau \cdot B_S}{\min(|S_t(x) - \gamma|, |S_{t-\tau}(x) - \gamma|)}.\]

For samples near the decision boundary (\(|S_t - \gamma|\) small), the
error can be \(O(1)\) even for modest \(\tau\).

\emph{Proof sketch.} By Lipschitz continuity (C3),
\(\|S_t - S_{t-\tau}\|_\infty \leq L_S \cdot \tau \cdot \beta_{\max} \cdot B_S\).
The probability of sign change is bounded by the probability that the
drift exceeds the margin. For small margin, this probability approaches
1. \(\square\)

\textbf{Status: CONJECTURED.} The Lipschitz bound is valid, but the
probability estimate requires a model of the drift direction, which
depends on the specific SCX update dynamics.

\subsubsection{2.6 Backlog-Induced Distribution
Shift}\label{backlog-induced-distribution-shift}

\textbf{Proposition 12.1 (Backlog Distribution Bias --- CONJECTURED).}
When samples are scored with delay \(\tau\), the effective acceptance
distribution is:

\[\tilde{P}_{S_t}(x,y) \propto \mathbf{1}\{S_{t-\tau(x)}(x,y) \geq \gamma\} \cdot P_0(x,y).\]

This distribution is biased toward samples where \(S_{t-\tau}\)
overestimates reliability relative to \(S_t\), because: - Samples scored
by an older (less accurate) gatekeeper are more likely to be
misclassified. - The direction of misclassification depends on whether
\(S_t\) has increased or decreased relative to \(S_{t-\tau}\).

\textbf{If \(S_t\) is monotonically improving}, then
\(S_t > S_{t-\tau}\) for most \((x,y)\) (the gatekeeper becomes more
generous over time). In this case, the stale gatekeeper is \textbf{more
conservative}, rejecting some samples that \(S_t\) would accept. The
effect is \textbf{reduced memory growth}, not incorrect acceptance.

\textbf{If \(S_t\) oscillates}, then \(S_{t-\tau}\) may be higher than
\(S_t\) in some regions, admitting samples that \(S_t\) would reject.
This introduces \textbf{noisy samples} into \(M_t\), degrading all
downstream estimates.

\subsubsection{2.7 Approximate Scoring for
Throughput}\label{approximate-scoring-for-throughput}

To prevent backlog, the gatekeeper can use \textbf{approximate scoring}:

\[\tilde{S}_t(x,y) = S_t(\text{NN}_k(x), y),\]

where \(\text{NN}_k(x)\) returns the \(k\) nearest neighbors of \(x\)
among already-scored samples.

\textbf{Theorem 12.6 (Approximation Error of NN Scoring --- PROVEN).}
Under C3 (Lipschitz gatekeeper) and assuming the already-scored samples
form an \(\varepsilon\)-net of the input space:

\[\sup_{x,y} |\tilde{S}_t(x,y) - S_t(x,y)| \leq L_S \cdot \varepsilon.\]

\emph{Proof.} By Lipschitz continuity,
\(|S_t(x) - S_t(\text{NN}(x))| \leq L_S \cdot \|x - \text{NN}(x)\| \leq L_S \cdot \varepsilon\)
when the scored samples form an \(\varepsilon\)-net. \(\square\)

\textbf{Status: PROVEN.} This provides a principled approximation with
controllable error.

\subsubsection{2.8 Mitigation Strategies}\label{mitigation-strategies-1}

\begin{longtable}[]{@{}
  >{\raggedright\arraybackslash}p{(\linewidth - 4\tabcolsep) * \real{0.3448}}
  >{\raggedright\arraybackslash}p{(\linewidth - 4\tabcolsep) * \real{0.2759}}
  >{\raggedright\arraybackslash}p{(\linewidth - 4\tabcolsep) * \real{0.3793}}@{}}
\toprule\noalign{}
\begin{minipage}[b]{\linewidth}\raggedright
Strategy
\end{minipage} & \begin{minipage}[b]{\linewidth}\raggedright
Effect
\end{minipage} & \begin{minipage}[b]{\linewidth}\raggedright
Trade-off
\end{minipage} \\
\midrule\noalign{}
\endhead
\bottomrule\noalign{}
\endlastfoot
\textbf{NN approximate scoring} & \(O(1)\) per sample (after index
build) & \(L_S \cdot \varepsilon\) score error \\
\textbf{Batch scoring} & Amortizes overhead & Increases delay for early
samples in batch \\
\textbf{Priority queue} & Score high-uncertainty samples first &
Requires uncertainty estimate before scoring (circular) \\
\textbf{Parallel/distributed experts} & Increases \(r_{\text{score}}\)
linearly with workers & Communication overhead \\
\textbf{Gatekeeper distillation} & Train a fast student gatekeeper
\(\tilde{S}\) to approximate \(S_t\) & Distillation error; requires
periodic retraining \\
\textbf{Adaptive sampling} & Reduce \(r_{\text{in}}\) when backlog
exceeds threshold & May miss novel samples \\
\end{longtable}

\begin{center}\rule{0.5\linewidth}{0.5pt}\end{center}

\subsection{3. Failure Mode 3: Calibration Breakdown (Client
Divergence)}\label{failure-mode-3-calibration-breakdown-client-divergence}

\subsubsection{3.1 Phenomenon}\label{phenomenon-2}

When two independent SCX clients (with different initial gatekeepers
\(S_0^{(1)}\), \(S_0^{(2)}\) or different data streams) evolve their
gatekeepers independently, their scoring functions may \textbf{diverge
irreconcilably}: they assign opposite reliability judgments to the same
samples, with no mechanism to reconcile.

\subsubsection{3.2 Formal Setup}\label{formal-setup}

\textbf{Definition 3.1 (Client Divergence).} Two clients \(A\) and \(B\)
evolve gatekeepers \(S_t^{(A)}\) and \(S_t^{(B)}\) from initial
conditions \(S_0^{(A)} \neq S_0^{(B)}\). Their \textbf{disagreement} at
time \(t\) is:

\[D_t(A, B) = \mathbb{P}_{(x,y) \sim P_0}\bigl(\text{sign}(S_t^{(A)}(x,y) - \gamma) \neq \text{sign}(S_t^{(B)}(x,y) - \gamma)\bigr).\]

Divergence is \textbf{irreconcilable} if
\(\liminf_{t \to \infty} D_t(A, B) > 0\).

\subsubsection{3.3 Existence of Multiple Fixed
Points}\label{existence-of-multiple-fixed-points}

\textbf{Theorem 12.7 (Multiple Fixed Points --- PROVEN).} The SCX
self-evolution dynamics can have multiple fixed points when: 1. The loss
landscape is non-convex (multiple local minima for the student). 2. The
gatekeeper's acceptance policy can be self-consistent at different
selectivity levels.

\emph{Proof by construction.} Consider two fixed points: - \textbf{Fixed
point A (conservative)}: \(S_A^*(x,y) \approx 0.3\) for all \((x,y)\).
The gatekeeper is conservative, rejecting most samples. The memory bank
is small, containing only the highest-confidence samples. The student is
trained on this small, clean subset. \textbf{Self-consistent}: the
conservative gatekeeper only admits samples it scores above 0.3, and the
student trained on those samples predicts labels that the gatekeeper is
comfortable with.

\begin{itemize}
\tightlist
\item
  \textbf{Fixed point B (permissive)}: \(S_B^*(x,y) \approx 0.7\) for
  all \((x,y)\). The gatekeeper is permissive, accepting most samples.
  The memory bank is large and diverse, including some noisy samples.
  The student is trained on this larger, noisier dataset.
  \textbf{Self-consistent}: the permissive gatekeeper admits most
  samples, and the student trained on those diverse samples makes
  predictions that maintain the permissive gatekeeper's calibration.
\end{itemize}

Both are valid fixed points of the dynamics. They are \textbf{not}
equivalent --- they disagree on approximately 40\% of decisions (samples
with true reliability between 0.3 and 0.7). \(\square\)

\textbf{Status: PROVEN.} The existence of multiple fixed points follows
from the non-convexity of the coupled system.

\subsubsection{3.4 Basin of Attraction}\label{basin-of-attraction}

\textbf{Definition 3.2 (Basin of Attraction).} The basin of attraction
\(\mathcal{B}(S^*, \theta^*)\) is the set of initial conditions
\((S_0, \theta_0)\) from which the system converges to
\((S^*, \theta^*)\).

\textbf{Conjecture 12.1 (Basin Structure).} The state space partitions
into basins of attraction:

\[\mathcal{Z} = \bigcup_{i=1}^{N_{\text{fp}}} \mathcal{B}_i \cup \mathcal{B}_{\text{boundary}},\]

where \(\mathcal{B}_i\) are the basins of the \(N_{\text{fp}}\) fixed
points, and \(\mathcal{B}_{\text{boundary}}\) is a measure-zero set of
initial conditions on the boundaries.

\emph{Status: \textbf{CONJECTURED.} The basin structure for coupled SA
systems is an active research area (Borkar, 2008, Chapter 9). The
conjecture is plausible but unproven for the specific SCX dynamics.}

\subsubsection{3.5 Condition for Irreconcilable
Divergence}\label{condition-for-irreconcilable-divergence}

\textbf{Theorem 12.8 (Irreconcilable Divergence Condition --- PROVEN).}
Two clients with initial conditions
\((S_0^{(A)}, \theta_0^{(A)}) \in \mathcal{B}_i\) and
\((S_0^{(B)}, \theta_0^{(B)}) \in \mathcal{B}_j\) with \(i \neq j\) will
converge to different fixed points with:

\[\lim_{t \to \infty} D_t(A, B) = D_{ij} > 0,\]

where \(D_{ij}\) is the disagreement rate between fixed points \(i\) and
\(j\).

\emph{Proof.} By definition of basins of attraction, client A converges
to fixed point \(i\) and client B to fixed point \(j\). Since
\(i \neq j\), there exists at least one \((x,y)\) where the fixed-point
gatekeepers disagree. By continuity of the scoring function, the
disagreement extends to a neighborhood of positive measure, so
\(D_{ij} > 0\). \(\square\)

\textbf{Status: PROVEN (conditional on basin structure).} The logic is
sound; the only gap is the rigorous characterization of basin boundaries
(Conjecture 12.1).

\subsubsection{3.6 Calibration Degradation
Rate}\label{calibration-degradation-rate}

\textbf{Proposition 12.2 (Calibration Drift Under Isolation ---
CONJECTURED).} If two clients start from \(\varepsilon\)-close initial
conditions but evolve on different data streams, their gatekeepers
diverge at rate:

\[\mathbb{E}[D_t(A, B)] \leq \varepsilon \cdot \exp(\lambda_{\max} \cdot t),\]

where \(\lambda_{\max}\) is the largest Lyapunov exponent of the
deterministic part of the dynamics.

\emph{Proof sketch.} Linearize the dynamics around the common
trajectory. The difference \(\delta_t = S_t^{(A)} - S_t^{(B)}\) evolves
as \(\delta_{t+1} \approx (I + \beta_t J_t) \delta_t + \text{noise}\),
where \(J_t\) is the Jacobian of SCXUpdate. The largest eigenvalue of
\(J_t\) determines the exponential divergence rate. \(\square\)

\textbf{Status: CONJECTURED.} The linearized analysis is standard for
dynamical systems (Strogatz, 2018), but deriving \(\lambda_{\max}\) for
the specific SCX dynamics requires further work.

\subsubsection{3.7 Mitigation Strategies}\label{mitigation-strategies-2}

\begin{longtable}[]{@{}
  >{\raggedright\arraybackslash}p{(\linewidth - 4\tabcolsep) * \real{0.3704}}
  >{\raggedright\arraybackslash}p{(\linewidth - 4\tabcolsep) * \real{0.4074}}
  >{\raggedright\arraybackslash}p{(\linewidth - 4\tabcolsep) * \real{0.2222}}@{}}
\toprule\noalign{}
\begin{minipage}[b]{\linewidth}\raggedright
Strategy
\end{minipage} & \begin{minipage}[b]{\linewidth}\raggedright
Mechanism
\end{minipage} & \begin{minipage}[b]{\linewidth}\raggedright
Cost
\end{minipage} \\
\midrule\noalign{}
\endhead
\bottomrule\noalign{}
\endlastfoot
\textbf{Federated gatekeeper averaging} & Periodically average
\(S_t^{(A)}\) and \(S_t^{(B)}\) & Requires communication; may pull both
toward a worse compromise \\
\textbf{Shared reference set} & Both clients evaluate against the same
\(M_0\) & Requires agreement on ground-truth labels for \(V_0\) \\
\textbf{Consensus anchoring} & Fix \(\hat{C}(x)\) as a shared anchor;
only evolve gatekeeper toward consensus, not away & Limits adaptation to
client-specific data distributions \\
\textbf{Cross-validation between clients} & Each client's memory bank
serves as validation for the other & Requires trust and data sharing \\
\textbf{Divergence detection + reset} & When
\(D_t(A,B) > \delta_{\text{threshold}}\), reset both to a common
initialization & Loses learned information \\
\end{longtable}

\begin{center}\rule{0.5\linewidth}{0.5pt}\end{center}

\subsection{4. Failure Mode 4: Adversarial Poisoning (Gatekeeper
Corruption)}\label{failure-mode-4-adversarial-poisoning-gatekeeper-corruption}

\subsubsection{4.1 Phenomenon}\label{phenomenon-3}

An adversary injects carefully crafted samples into the data stream,
causing the gatekeeper to systematically misclassify clean samples as
noise (or vice versa). The self-evolution loop amplifies the corruption
by training the student on poisoned data.

\subsubsection{4.2 Threat Model}\label{threat-model}

\textbf{Definition 4.1 (Adversary Capabilities).} The adversary can: -
\textbf{Injection}: Add up to \(\eta_{\text{adv}}\) fraction of samples
to the incoming data stream. - \textbf{Knowledge}: Know the current
gatekeeper \(S_t\), the expert models \(\{f_m\}\), and the feature
representation \(\phi\). - \textbf{Objective}: Maximize the false
positive rate (clean → noise) or false negative rate (noise → clean) of
the converged gatekeeper.

We assume the adversary does \textbf{not} control the expert models or
the NEP training procedure --- only the data fed to the gatekeeper.

\subsubsection{4.3 Adversarial Sample
Construction}\label{adversarial-sample-construction}

\textbf{Definition 4.2 (Consensus-Minimizing Adversarial Sample).} To
cause the gatekeeper to reject a clean sample \((x, y_{\text{true}})\),
the adversary constructs:

\[(x_{\text{adv}}, y_{\text{adv}}) = \arg\min_{(x',y') \in \mathcal{B}_\varepsilon(x)} \hat{C}(x') \quad \text{subject to} \quad \ell(f_m(x'), y') > \tau \text{ for most } m,\]

i.e., an \(\varepsilon\)-perturbation of a clean sample that maximally
confuses the expert ensemble while remaining indistinguishable to the
gatekeeper's feature representation.

\subsubsection{4.4 Poisoning Amplification via
Self-Evolution}\label{poisoning-amplification-via-self-evolution}

The key danger in SCX self-evolution is \textbf{amplification}: one
round of poisoning corrupts the gatekeeper, which admits more poisoned
samples, which further corrupts the gatekeeper.

\textbf{Theorem 12.9 (Poisoning Amplification Factor --- CONJECTURED).}
If the adversary injects a fraction \(\eta_{\text{adv}}\) of poisoned
samples at each round, the effective corruption after \(T\) rounds is:

\[\eta_{\text{eff}}(T) \geq \eta_{\text{adv}} \cdot \frac{1 - \gamma^T}{1 - \gamma},\]

where
\(\gamma = \mathbb{P}_{P_0}(S_t(x_{\text{adv}}, y_{\text{adv}}) \geq \gamma_t \mid \text{poisoned})\)
is the probability that poisoned samples pass the gatekeeper.

When \(\gamma > 0.5\) (poisoned samples are more likely to be accepted
than rejected), the amplification diverges geometrically:

\[\eta_{\text{eff}}(T) \to 1 \quad \text{as} \quad T \to \infty.\]

\emph{Proof sketch.} At each round, a fraction of poisoned samples enter
\(M_t\), corrupting the student. The corrupted student provides feedback
that shifts \(S_{t+1}\) toward accepting more poisoned samples (because
the student agrees with the poisoned labels). This increases \(\gamma\)
over time, creating a positive feedback loop. \(\square\)

\textbf{Status: CONJECTURED.} The geometric amplification depends on
\(\gamma > 0.5\), which in turn depends on the adversary's ability to
craft samples that both confuse experts and appear reliable to the
current gatekeeper. This is a strong assumption that may not hold for
well-regularized gatekeepers.

\subsubsection{4.5 Detection Boundary for Adversarial
Samples}\label{detection-boundary-for-adversarial-samples}

\textbf{Theorem 12.10 (Adversarial Detection Boundary --- PARTIALLY
PROVEN).} Under the consensus score \(\hat{C}(x)\) computed from \(M\)
i.i.d. experts, an adversarial perturbation \(\delta\) on a clean sample
\(x\) changes the expected consensus by at most:

\[\mathbb{E}[|\hat{C}(x + \delta) - \hat{C}(x)|] \leq \frac{1}{M} \sum_{m=1}^M \mathbb{P}(\ell(f_m(x+\delta), y) > \tau \neq \ell(f_m(x), y) > \tau).\]

If each expert has a \textbf{robustness radius} \(\rho_m\) such that
\(\|\delta\| \leq \rho_m \implies f_m(x+\delta) = f_m(x)\), then for
\(\|\delta\| \leq \min_m \rho_m\), the consensus score is
\textbf{unchanged} and the adversarial sample is undetectable by
consensus alone.

\emph{Proof.} The consensus changes only if the adversary's perturbation
flips at least one expert's error indicator. If all experts are robust
within radius \(\rho_{\min}\), then no expert error indicator flips, and
\(\hat{C}\) is unchanged. \(\square\)

\textbf{Status: PROVEN.} This follows from the definition of \(\hat{C}\)
and provides a clean condition for undetectability.

\subsubsection{4.6 Gatekeeper Robustness
Certificate}\label{gatekeeper-robustness-certificate}

\textbf{Proposition 12.3 (Lipschitz-Based Robustness --- PROVEN).} Under
C3 (Lipschitz gatekeeper with constant \(L_S\)), for any perturbation
\(\delta\):

\[|S_t(x + \delta, y) - S_t(x, y)| \leq L_S \cdot \|\phi(x + \delta) - \phi(x)\| \leq L_S \cdot L_\phi \cdot \|\delta\|,\]

where \(L_\phi\) is the Lipschitz constant of the feature map. The
gatekeeper is \textbf{certifiably robust} for perturbations with
\(\|\delta\| \leq \frac{|S_t(x,y) - \gamma|}{L_S \cdot L_\phi}\).

\emph{Proof.} Direct from C3 and the Lipschitz property of \(\phi\).
\(\square\)

\textbf{Status: PROVEN.} This provides a certified radius within which
adversarial perturbations cannot change the gatekeeper's decision.

\subsubsection{4.7 When the Student is the
Vulnerability}\label{when-the-student-is-the-vulnerability}

Even if the gatekeeper is robust, the \textbf{NEP student} may be
vulnerable. If poisoned samples enter \(M_t\): 1. The student is trained
on poisoned labels. 2. The student's predictions on clean samples become
corrupted. 3. The corrupted student feedback shifts the gatekeeper. 4.
The cycle continues.

\textbf{Theorem 12.11 (Student Poisoning Vulnerability ---
CONJECTURED).} Under linear regression with \(\ell_2\) loss, if an
\(\eta_{\text{adv}}\) fraction of training labels are flipped, the
parameter error is:

\[\mathbb{E}[\|\hat{\theta} - \theta^*\|^2] \geq \eta_{\text{adv}}^2 \cdot \frac{\mathbb{E}[\|x\|^2]}{\lambda_{\min}(\Sigma)^2},\]

where \(\Sigma = \mathbb{E}[x x^\top]\) is the feature covariance. For
neural networks, the error may be larger due to the non-linear
amplification of poisoned gradients.

\emph{Status: \textbf{CONJECTURED} (for neural networks). The linear
case is standard (Hampel et al., 1986).}

\subsubsection{4.8 Mitigation Strategies}\label{mitigation-strategies-3}

\begin{longtable}[]{@{}
  >{\raggedright\arraybackslash}p{(\linewidth - 4\tabcolsep) * \real{0.3030}}
  >{\raggedright\arraybackslash}p{(\linewidth - 4\tabcolsep) * \real{0.3333}}
  >{\raggedright\arraybackslash}p{(\linewidth - 4\tabcolsep) * \real{0.3636}}@{}}
\toprule\noalign{}
\begin{minipage}[b]{\linewidth}\raggedright
Strategy
\end{minipage} & \begin{minipage}[b]{\linewidth}\raggedright
Mechanism
\end{minipage} & \begin{minipage}[b]{\linewidth}\raggedright
Limitation
\end{minipage} \\
\midrule\noalign{}
\endhead
\bottomrule\noalign{}
\endlastfoot
\textbf{Robust aggregation (median-of-experts)} & Replace mean consensus
\(\hat{C}\) with median or trimmed mean & Reduces effective \(M\);
requires \(>50\%\) clean experts \\
\textbf{Anomaly detection on \(\phi(x)\)} & Flag samples with atypical
features as potential adversarial & Adversary can constrain
\(\|\delta\|\) to stay in-distribution \\
\textbf{Certified robustness (Lipschitz)} & Bound
\(|S_t(x+\delta) - S_t(x)| \leq L_S L_\phi \|\delta\|\) & Loose bounds
for large \(L_S, L_\phi\) \\
\textbf{Differential privacy in gatekeeper update} & Add noise to
\(\Delta S_t\) to mask individual sample influence & Reduces convergence
rate by \(O(\sigma_{\text{DP}}^2)\) \\
\textbf{Student robust training} & Train NEP with adversarial data
augmentation & Increases training cost; may reduce clean accuracy \\
\textbf{Outlier-resistant loss} & Use Huber or Tukey loss for student
training & Less statistically efficient than \(\ell_2\) for clean
data \\
\textbf{Human-in-the-loop for high-stakes decisions} & Flag samples
where \(\hat{C}(x)\) and \(S_t(x)\) disagree strongly & Does not
scale \\
\end{longtable}

\begin{center}\rule{0.5\linewidth}{0.5pt}\end{center}

\subsection{5. Interaction Effects: Compound
Failures}\label{interaction-effects-compound-failures}

\subsubsection{5.1 Backlog + Premature
Freezing}\label{backlog-premature-freezing}

When backlog (Section 2) \textbf{and} premature freezing (Section 1)
co-occur: - The backlog means the gatekeeper's decisions are stale
(\(S_{t-\tau}\) instead of \(S_t\)). - Premature freezing means the
gatekeeper stops updating even as backlog grows. - \textbf{Compound
effect}: The effective gatekeeper \(S_{t-\tau}\) freezes at an even
earlier (worse) state, and the delay \(\tau\) grows without bound. The
system becomes \textbf{doubly frozen} --- neither the gatekeeper nor the
scoring process can recover.

\subsubsection{5.2 Adversarial Poisoning + Client
Divergence}\label{adversarial-poisoning-client-divergence}

When adversarial poisoning (Section 4) targets one of two clients: -
Client A (poisoned) converges to a corrupted fixed point. - Client B
(clean) converges to a clean fixed point. - Their disagreement
\(D_t(A, B)\) grows over time. - \textbf{Compound effect}: The adversary
can \textbf{maximize} \(D_t(A, B)\) by injecting samples that are
specifically chosen to push A's gatekeeper away from B's. This is a
\textbf{targeted divergence attack}: the adversary exploits the
multi-fixed-point structure.

\subsubsection{5.3 Backlog + Adversarial
Poisoning}\label{backlog-adversarial-poisoning}

When backlog means samples are scored by stale gatekeeper
\(S_{t-\tau}\): - The adversary can inject samples that are
\textbf{benign} to the current gatekeeper \(S_t\) but \textbf{malignant}
to the stale gatekeeper \(S_{t-\tau}\). - Since scoring uses
\(S_{t-\tau}\), the malignant samples pass through. - \textbf{Compound
effect}: Backlog creates a \textbf{temporal attack surface} --- the
adversary can exploit the gap between the current and stale gatekeepers.

\begin{center}\rule{0.5\linewidth}{0.5pt}\end{center}

\subsection{6. Diagnostic Tests for Each Failure
Mode}\label{diagnostic-tests-for-each-failure-mode}

\subsubsection{6.1 Premature Freezing
Detection}\label{premature-freezing-detection}

\textbf{Test 1 (Gradient Norm Monitor).} Track \(\|\Delta S_t\|_{M_0}\)
over a sliding window of \(W\) steps. If
\(\max_{i \in [t-W, t]} \|\Delta S_i\| < \varepsilon_{\text{diag}}\) for
\(W\) consecutive windows, flag as frozen.

\textbf{Test 2 (Student Improvement Correlation).} Track the correlation
between \(\|\Delta S_t\|\) and \(\Delta \Psi_{\text{student}, t}\). If
the gatekeeper is frozen but the student is still improving, the
correlation drops to zero.

\subsubsection{6.2 Backlog Detection}\label{backlog-detection}

\textbf{Test 3 (Scoring Delay Monitor).} Track
\(\tau(x) = t_{\text{current}} - t_{\text{scored}}\) for each sample. If
\(\text{median}(\tau) > \tau_{\text{threshold}}\), flag backlog.

\textbf{Test 4 (Memory Growth vs.~Scoring Rate).} Track \(dN_t/dt\)
(memory growth rate) and \(d\text{Scored}_t/dt\) (scoring rate). If the
former consistently exceeds the latter, backlog is growing.

\subsubsection{6.3 Calibration Divergence
Detection}\label{calibration-divergence-detection}

\textbf{Test 5 (Inter-Client Disagreement).} Periodically compare
gatekeeper decisions between clients on a shared validation set. If
\(D_t(A, B)\) is increasing and exceeds \(\delta_{\text{threshold}}\),
flag divergence.

\textbf{Test 6 (Fixed-Point Distance Estimate).} Estimate the distance
between \(S_t^{(A)}\) and \(S_t^{(B)}\) in function space:
\(\|S_t^{(A)} - S_t^{(B)}\|_{M_0}\). A monotonically increasing trend
suggests divergence to different fixed points.

\subsubsection{6.4 Adversarial Poisoning
Detection}\label{adversarial-poisoning-detection}

\textbf{Test 7 (Consensus-Gatekeeper Discrepancy).} Track
\(\mathbb{E}[|\hat{C}(x) - S_t(x, y(x))|]\) on incoming samples. A
sudden increase may indicate adversarial samples that confuse experts
but are scored highly by the gatekeeper.

\textbf{Test 8 (Expert Agreement Anomaly).} Track the per-expert
agreement matrix. Adversarial samples may cause unusual patterns (e.g.,
all experts agree on the wrong answer, or expert disagreement is bimodal
rather than uniform).

\textbf{Test 9 (Student Loss on Held-Out Clean Set).} Monitor
\(\Psi_{\text{student}}\) on \(V_0\). Poisoning should cause this to
\textbf{increase} even as the training loss on \(M_t\) decreases (a
classic sign of label poisoning).

\begin{center}\rule{0.5\linewidth}{0.5pt}\end{center}

\subsection{7. Summary Table}\label{summary-table}

\begin{longtable}[]{@{}
  >{\raggedright\arraybackslash}p{(\linewidth - 10\tabcolsep) * \real{0.1857}}
  >{\raggedright\arraybackslash}p{(\linewidth - 10\tabcolsep) * \real{0.1000}}
  >{\raggedright\arraybackslash}p{(\linewidth - 10\tabcolsep) * \real{0.1143}}
  >{\raggedright\arraybackslash}p{(\linewidth - 10\tabcolsep) * \real{0.2714}}
  >{\raggedright\arraybackslash}p{(\linewidth - 10\tabcolsep) * \real{0.1429}}
  >{\raggedright\arraybackslash}p{(\linewidth - 10\tabcolsep) * \real{0.1857}}@{}}
\toprule\noalign{}
\begin{minipage}[b]{\linewidth}\raggedright
Failure Mode
\end{minipage} & \begin{minipage}[b]{\linewidth}\raggedright
Cause
\end{minipage} & \begin{minipage}[b]{\linewidth}\raggedright
Effect
\end{minipage} & \begin{minipage}[b]{\linewidth}\raggedright
Critical Parameter
\end{minipage} & \begin{minipage}[b]{\linewidth}\raggedright
Severity
\end{minipage} & \begin{minipage}[b]{\linewidth}\raggedright
Provability
\end{minipage} \\
\midrule\noalign{}
\endhead
\bottomrule\noalign{}
\endlastfoot
\textbf{1. Premature Freezing} & \(\beta_t\) decays too fast relative to
student convergence & Gatekeeper trapped at suboptimal
\(S^*_{\text{sub}}\); student converges to wrong target &
\(t_{\text{freeze}} = (\beta_0 B_S / \varepsilon_{\text{mach}})^{1/b}\)
& \textbf{MODERATE} (rare with double precision) & Condition:
\textbf{PROVEN}; gap: \textbf{CONJECTURED} \\
\textbf{2. Backlog} & \(r_{\text{in}} > r_{\text{score}}\) & Stale
gatekeeper decisions; memory bank contaminated with misclassified
samples & \(r_{\text{in}} / r_{\text{score}}\) & \textbf{HIGH}
(practical for large \(M\), high data rate) & Condition:
\textbf{PROVEN}; effect: \textbf{CONJECTURED} \\
\textbf{3. Client Divergence} & Multiple fixed points; different initial
conditions or data streams & Irreconcilable gatekeeper disagreement; no
convergence to consensus & \(D_{ij}\) (inter-fixed-point disagreement) &
\textbf{MODERATE} (only in multi-client deployment) & Existence:
\textbf{PROVEN}; basin structure: \textbf{CONJECTURED} \\
\textbf{4. Adversarial Poisoning} & Adversary injects crafted samples &
Gatekeeper corruption amplifies through self-evolution loop &
\(\eta_{\text{adv}}\), \(\gamma\) (acceptance rate of poisoned samples)
& \textbf{HIGH} (severe with amplification) & Amplification:
\textbf{CONJECTURED}; robustness certificate: \textbf{PROVEN} \\
\end{longtable}

\subsubsection{7.1 Risk Matrix}\label{risk-matrix}

\begin{verbatim}
                      Likelihood
                    Low     Med    High
Severity  High     |       |  2   |  4   |
          Moderate |   1   |  3   |      |
          Low      |       |      |      |
\end{verbatim}

\begin{itemize}
\tightlist
\item
  \textbf{High risk, high likelihood}: Backlog (2), Adversarial
  poisoning (4)
\item
  \textbf{High risk, moderate likelihood}: ---
\item
  \textbf{Moderate risk, low likelihood}: Premature freezing (1)
\item
  \textbf{Moderate risk, moderate likelihood}: Client divergence (3)
\end{itemize}

\subsubsection{7.2 Design Recommendations}\label{design-recommendations}

\begin{enumerate}
\def\labelenumi{\arabic{enumi}.}
\item
  \textbf{Prevent premature freezing}: Always maintain
  \(\beta_{\min} > 0\) or use cyclical learning rates. Monitor
  \(\|\Delta S_t\|\) as a diagnostic.
\item
  \textbf{Prevent backlog}: Implement NN approximate scoring with
  bounded error. Monitor scoring delay. Trigger adaptive sampling when
  backlog exceeds threshold.
\item
  \textbf{Prevent client divergence}: Use a shared reference set \(M_0\)
  with agreed-upon labels for \(V_0\). Periodically synchronize
  gatekeepers in federated deployments.
\item
  \textbf{Prevent adversarial poisoning}: Use Lipschitz-regularized
  gatekeepers with certified robustness radii. Train the NEP student
  with robust loss functions. Monitor consensus-gatekeeper discrepancy
  on held-out data.
\end{enumerate}

\begin{center}\rule{0.5\linewidth}{0.5pt}\end{center}

\emph{End of 12\_edge\_cases.md --- Edge Cases and Failure Modes}

\begin{center}\rule{0.5\linewidth}{0.5pt}\end{center}

\subsection{References}\label{references}

\begin{enumerate}
\def\labelenumi{\arabic{enumi}.}
\tightlist
\item
  Strogatz, S. H. (2018). \emph{Nonlinear Dynamics and Chaos} (2nd ed.).
  CRC Press.
\item
  Borkar, V. S. (2008). \emph{Stochastic Approximation: A Dynamical
  Systems Viewpoint}. Cambridge University Press.
\item
  Hampel, F. R., et al.~(1986). \emph{Robust Statistics: The Approach
  Based on Influence Functions}. Wiley.
\item
  Goodfellow, I. J., Shlens, J., \& Szegedy, C. (2015). Explaining and
  harnessing adversarial examples. \emph{ICLR 2015}.
\item
  Madry, A., et al.~(2018). Towards deep learning models resistant to
  adversarial attacks. \emph{ICLR 2018}.
\item
  Biggio, B., Nelson, B., \& Laskov, P. (2012). Poisoning attacks
  against support vector machines. \emph{ICML 2012}.
\item
  Steinhardt, J., Koh, P. W., \& Liang, P. (2017). Certified defenses
  for data poisoning attacks. \emph{NIPS 2017}.
\item
  Blanchard, P., et al.~(2017). Machine learning with adversaries:
  Byzantine tolerant gradient descent. \emph{NIPS 2017}.
\item
  Kushner, H. J., \& Yin, G. G. (2003). \emph{Stochastic Approximation
  and Recursive Algorithms and Applications} (2nd ed.). Springer.
\item
  Polyak, B. T., \& Juditsky, A. B. (1992). Acceleration of stochastic
  approximation by averaging. \emph{SIAM Journal on Control and
  Optimization}, 30(4), 838-855.
\end{enumerate}

\begin{center}\rule{0.5\linewidth}{0.5pt}\end{center}

\subsection{Changelog}\label{changelog}

\begin{longtable}[]{@{}
  >{\raggedright\arraybackslash}p{(\linewidth - 6\tabcolsep) * \real{0.1875}}
  >{\raggedright\arraybackslash}p{(\linewidth - 6\tabcolsep) * \real{0.2500}}
  >{\raggedright\arraybackslash}p{(\linewidth - 6\tabcolsep) * \real{0.2500}}
  >{\raggedright\arraybackslash}p{(\linewidth - 6\tabcolsep) * \real{0.3125}}@{}}
\toprule\noalign{}
\begin{minipage}[b]{\linewidth}\raggedright
Date
\end{minipage} & \begin{minipage}[b]{\linewidth}\raggedright
Defect
\end{minipage} & \begin{minipage}[b]{\linewidth}\raggedright
Change
\end{minipage} & \begin{minipage}[b]{\linewidth}\raggedright
Severity
\end{minipage} \\
\midrule\noalign{}
\endhead
\bottomrule\noalign{}
\endlastfoot
2026-06-28 & B4 & \textbf{Renamed gatekeeper update rate from η to β}:
The symbol \(\eta\) is reserved throughout SCX theory for the global
noise rate \(\eta = P(Z=1)\). All local uses of \(\eta(t)\) as
gatekeeper update rate replaced with \(\beta_t\). Notation note added at
document header and Theorem 12.1. Fixes FAIL-2 from cross-reference
audit. & FATAL \\
\end{longtable}

\end{document}
